% Part: sets-functions-relations
% Chapter: relations
% Section: reflections
%
\documentclass[../../../include/open-logic-section]{subfiles}

\begin{document}

\olfileid{sfr}{rel}{ref}
\olsection{فلسفي غور}

په \olref[set]{sec} کښې مو اړيکې د خاصو سټونو په توګه تعريف کړې۔
لږ تم کېږو او يوه لنډه فلسفي پوښتنه کوو: داسې تعريف څۀ \emph{کوي}؟
دا ډېره د شک خبره ده چې مونږ دې ووايو چې د ماوراءالطبيعتي عينيت ځينې
حقيقتونه مو \emph{کشف} کړي دي؛ مثلاً، دا چې په $\Nat$ باندې ترتيبي
اړيکه \emph{په اصل کښې} هغه سټ
$R=  \Setabs{\tuple{n,m}}{n, m \in \Nat\text{ او } n < m}$
\emph{وختله} چې په \olref[set]{sec} کښې مو تعريف کړے و۔ د دې خبرې
درې وجهې دا دي۔

لومړۍ: په \olref[set][pai]{wienerkuratowski} کښې مو
$\tuple{a, b} = \{\{a\}, \{a, b\}\}$ تعريف کړے و۔ د دې پر ځاے
دا تعريف وګورئ:
$\lVert a, b\rVert = \{\{b\}, \{a, b\}\} = \tuple{b,a}$۔
کله چې $a \neq b$ وي، نو $\tuple{a, b} \neq \lVert a,b\rVert$۔
خو د $\tuple{a,b}$ پر ځاے مو همداسې $\lVert a,b\rVert$ هم د مرتبې
جوړې تعريف ګڼلے شو۔ دواړو تعريفونو به يو شان ښۀ کار کړے و۔ نو اوس
د طبيعي عددونو ترتيبي اړيکه ``کېدو'' ته دوه يو شان ښۀ نوماندان لرو:
\begin{align*}
		R &= \Setabs{\tuple{n,m}}{n, m \in \Nat \text{ او }n < m}\\
		S &= \Setabs{\lVert n,m\rVert}{n, m \in \Nat \text{ او }n < m}.
\end{align*}
ځکه چې د غړو له مخې د برابرۍ د اصل له مخې $R \neq S$ دي، ښکاره
ده چې \emph{دواړه} په $\Nat$ باندې له ترتيبي اړيکې سره عينيت نۀ شي
لرلے۔ خو دا به محض خپله خوښه، او له دې امله يو څه د شرم خبره وي، چې
ادعا وکړو چې په واقعيت کښې $R$، د $S$ پر ځاے (يا برعکس)، ترتيبي
اړيکه \emph{ده}۔ (دا د سټونو ته د راکمونې د نظر خلاف د هغه دليل يوه
ډېره ساده بېلګه ده چې بناسراف په \citeyear{Benacerraf1965} کښې مشهور
کړے و۔ مونږ به څو ځله بيا دې ته راګرځو۔)

دويمه: که فکر کوو چې \emph{هره} اړيکه بايد له يو سټ سره عينيت
ولري، نو پخپله د سټ د غړيتوب اړيکه، $\in$، هم بايد يو خاص سټ وي۔
په رښتيا، هغه به خامخا سټ $\Setabs{\tuple{x,y}}{x \in y}$ وي۔
خو آيا دا سټ شته؟ د رسل پاراډوکس ته په کتو، د داسې سټ د موجوديت
ادعا ساده خبره نۀ ده۔ په حقيقت کښې،
\oliflabeldef{cumul:::part}{د سټونو هغه نظريه چې مونږ ئې په
\olref[cumul][][]{part} کښې جوړوو، د دې سټ موجوديت به \emph{رد}
کړي۔\footnote{له مخه ئې وجه دا ده۔ د تناقض له لارې د ثبوت دپاره فرض
کړئ چې $I = \Setabs{\tuple{x,y}}{x \in y}$ شته۔ بيا
$\bigcup \bigcup I$ ټولشموله سټ دے، چې له
\olref[sfr][z][sep]{thm:NoUniversalSet} سره تناقض لري۔}}{د سټونو
نظريه په يوه دقيقه طريقه، د بديهي اصولو د يوې نظريې په توګه، جوړول
ممکن دي، او هغه نظريه به په رښتيا د دې سټ موجوديت رد کړي۔}
نو که ځينې اړيکې د سټونو په توګه هم ګڼلے شو، د سټ د غړيتوب اړيکه
به يوه خاصه پېښه وي۔

درېيمه: کله چې اړيکې له سټونو سره ``يوې'' ګڼو، ومو وئيل چې ځان ته
به د $\tuple{x,y} \in R$ پر ځاے د $Rxy$ ليکلو اجازه ورکوو۔ دا سمه
ده، په دې شرط چې د غړيتوب اړيکه، ``$\in$''، د يو پريديکات
\emph{په توګه} وکارول شي۔ خو که فکر کوو چې ``$\in$'' د سټ د يو
خاص ډول نښه ده، نو عبارت ``$\tuple{x,y} \in R$'' يوازې له درېيو
مفردو اصطلاحاتو جوړ دے چې سټونه ښيي: ``$\tuple{x,y}$''، ``$\in$''
او ``$R$''۔ د نومونو دا ډول لړ، لکه بې معنا توريز تسلسل ``پياله
قلمدان مېز''، د يو بیان د څرګندولو وس نۀ لري۔ نو بيا هم، که ځينې
اړيکې د سټونو په توګه ګڼلے شو، د سټ د غړيتوب اړيکه بايد يوه خاصه
پېښه وي۔ (په دې کښې د فرېګه د \emph{آس} د مفهوم د پاراډوکس يوه
ساده بڼه او هغه مشهور اعتراض سره يوځاے شوي چې وېټګنشټاين يو وخت
پر رسل کړے و۔)

نو دې خبرې مونږ کوم ځاے ته ورسولو؟ دا وئيل چې په عددونو باندې
اړيکې سټونه دي، هېڅ \emph{غلط} نۀ دي۔ يوازې بايد په هغې موخې پوه
شو چې دا خبره ورسره کوو۔ مونږ د ماوراءالطبيعتي عينيت يو حقيقت نۀ
بيانوو۔ يوازې دا يادونه کوو چې په ځينو سياقونو کښې (ځينې) اړيکې
خاص سټونه \emph{ګڼلے} شو (او ګڼو به ئې)۔

\end{document}
